% !TeX root = ../tfg.tex
% !TeX encoding = utf8
%
% ====================================================================
%  APÉNDICE A — PREGUNTAS REALIZADAS A LA EMPRESA Y RESPUESTAS
%
%  RECORDATORIO AL EQUIPO:
%  - Este apéndice es OBLIGATORIO según las instrucciones de la asignatura.
%  - Debéis incluir las preguntas que realizasteis en la entrevista/visita
%    a la empresa y las respuestas que os dieron.
%  - Organizad las preguntas por temática (Tema 2 y Tema 3) o por el
%    orden en que se realizaron.
%  - Podéis incluir el nombre o cargo del entrevistado (si da permiso)
%    y la fecha y medio de la entrevista (presencial, videollamada...).
%  - Si hay información que el entrevistado pidió que no se publicara,
%    indicadlo (p.ej. "Dato confidencial — no se puede reproducir").
%  - Si grabasteis la entrevista (con permiso), podéis transcribir los
%    fragmentos más relevantes.
% ====================================================================

\chapter{Entrevista a CaixaBank: Temas 2 y 3 (análisis de puestos y planificación)}

% ============================================================
\section{Datos de la entrevista}
% ============================================================

\begin{table}[H]
  \centering
  \begin{tabularx}{\textwidth}{lX}
    \toprule
    \textbf{Parámetro}       & \textbf{Detalle} \\
    \midrule
    Empresa                  & CaixaBank \\
    Nombre del entrevistado  & [Anónimo a petición] \\
    Cargo                    & [Director de Recursos Humanos regional] \\
    Departamento             & [Recursos Humanos] \\
    Fecha                    & [13/03/2026] \\
    Duración                 & [15 min] \\
    Formato                  & [Telefónica] \\
    \bottomrule
  \end{tabularx}
\end{table}

% ============================================================
\section{Preguntas sobre el Tema 2 — Análisis y diseño de puestos}
% ============================================================

%Jose Angel
\subsubsection*{En clase hemos estudiado que el Análisis del Puesto de Trabajo es el proceso a través del cual la empresa recopila y analiza la información sobre
 el puesto para identificar sus tareas y responsabilidades. En una entidad tan grande como CaixaBank, cuando necesitáis definir un puesto
(por ejemplo, de nueva creación o uno que ha cambiado mucho), ¿cómo es vuestro proceso paso a paso para analizarlo? ¿Sentís que este proceso os funciona
 bien o es demasiado burocrático?}

\textbf{Respuesta:}\newline
``Los puestos los define un departamento Secretaria Tecnica de Negocios. Deloitte desarrollo el modelo. De abajo arriba, desde las necesidades del cliente. No se hace desde recursos humanos la definición, esta externalizado. Se basan mucho en las soft skills (dependiendo del lugar donde está la oficina).''
%David

\subsubsection*{Sabemos que existen algunos trabajos que pueden resultar monotonos o terminar quemando al empleado. ¿Qué medidas tomáis para rediseñar puestos monótonos o, por el contrario, muy especializados? ¿Habéis aplicado técnicas como la rotación de puestos, la ampliación de tareas o dar más autonomía a los empleados?}

\textbf{Respuesta:}\newline
``Desde la empresa se trabaja desde dos ámbitos: Una forma reactiva, de forma que se publican las vacantes frecuentemente y cada trabajador puede postular al que este interesado; y otra forma proactiva, en la que se hacen entrevistas individuales para ver si un empleado concreto encaja con el puesto para el cual estamos buscando empleado.''

\subsubsection*{En CaixaBank, a la hora de definir lo que necesitáis de un trabajador, ¿os centráis más en una lista estricta de tareas o vais más hacia un modelo de competencias blandas y adaptabilidad?}

\textbf{Respuesta:}\newline
``Los puestos se definen principalmente basandose en las tareas que van er, de forma que nos centramos principalmente en las "Hard-skills. Existe un departamento técnico específico que define las necesidades de cada puesto.''

% ============================================================
\section{Preguntas sobre el Tema 3 — Planificación de RRHH}
% ============================================================

% JESÚS 
\subsubsection*{Pasando a un plano más estratégico, la planificación de recursos humanos trata de anticiparse y prever el movimiento de personas. ¿Cómo planificáis las necesidades de personal a futuro? ¿Dirías que es una mera planificación administrativa de cuánta gente falta o está totalmente alineada con la estrategia de negocio del banco, es decir, pensáis en las necesidades a corto plazo (jubilaciones, etc.) o lo planificáis más a largo plazo?}

\textbf{Respuesta:}\newline
``En el tema de dimensionamineto se tienen en cuenta muchísimos factores y es un aspecto clave al determinar toda la estructura de costes. Los factores más importante son el mercado actual y potencial de cada plaza (vacante o puesto), la su estacionalidad o la tipología de clientes (por ejemplo, suelen tener distintas necesidades en pueblos que en ciudades) entre otros que son fundamentales para algo muy importante para nosotros como es el dimensionamiento.''

\subsubsection*{¿Y hay alguna herramienta o método concreto que useis para estas previones de demanda u oferta de trabajadores?}

\textbf{Respuesta:}\newline
``Se está trabajando con IA y, como la inmensa mayoría de las empresas, en todo el tema de planificación se usa SAP, que es un programa de gestión (SIRH) desde hace muchísimos años que ha evolucionado y que se usa en empresas desde el ámbito de la construcción hasta recursos humanos (casi todas las grandes empresas lo usan para eso). Respecto a la IA tenemos licencias de Copilot.''

% JULIAN ----------------------------------------------------------------------------------------------------
% - ¿Hay Sistemas de Información de RRHH? ¿En qué consisten? ¿Cómo pueden mejorar? 
% - Procesos Sustractivos: DEMANDA MENOR A OFERTA 
%   o ¿Existe suspensión laboral (excedencias/ERTES)? 
%   o ¿Existe ruptura laboral voluntaria (dimisión, jubilación, jubilación anticipada)? 

\subsubsection*{Para manejar a tantos miles de empleados, imaginamos que la tecnología es vital. Los sistemas de información de recursos humanos (SIRH) recopilan, registran y almacenan los datos de los empleados. ¿Qué tipo de plataforma utilizáis en CaixaBank para esto? ¿Estáis contentos con ella o creéis que le falta evolucionar en algún aspecto?”}

\textbf{Respuesta:}\newline
``Ahora mismo se esta trabajando con IA, y bueno, nosotros como la inmensa mayoría de empresas, utilizamos una herramienta llamada SAP que se utiliza en todos los procesos de planificación. SAP es un programa de gestión desde hace muchísimos años. Todas las grandes empresas trabajamos en el entorno SAP. Y ahora como te digo, también herramientas de IA, en concreto tenemos como herramienta de cabecera y licenciada Copilot. ''

\subsubsection*{Hablando del movimiento de personas, a veces los trabajadores pausan su actividad, interrumpiendo temporalmente la prestación laboral. ¿Cómo gestionáis temas como las excedencias voluntarias o por cuidado de familiares? ¿Es fácil reincorporar el talento después?}
\textbf{Respuesta:}\newline
`` Siempre que el empleado se preste (en la mayoría de los casos lo hace), realizamos entrevistas de salida y tratamos de organizar cual ha sido el punto de valor que ha provocado la ruptura. ''

\subsubsection*{Y cuando alguien decide irse definitivamente, ya sea por dimisión o por jubilación (normal o anticipada), ¿cómo gestionáis esas salidas? ¿Hacéis entrevistas de salida para entender los motivos y mejorar la retención del talento?}
\textbf{Respuesta:}\newline
`` Para la retención de talento, aunque a veces se toma por medidas económicas (donde tenemos una libertad limitada), intentamos no mantener exclusivamente por motivos económicos. Lo que trabajamos mucho es en el desarrollo profesional dentro de la empresa. Trabajamos mucho el tema de la formación, sobretodo en puestos técnicos donde hay una mayor tasa de abandono. Desarrollamos un plan de carrera en base a expectativas, y el tema de oportunidades, dado que nuestra empresa tiene una capacidad que muy pocas empresas en este país tienen.''

% ------------------PARTE DE ISMAEL-----------------

\subsubsection*{En ocasiones, por estrategia o economía, es la empresa quien debe finalizar la relación laboral. En un sector de tanta transformación como la banca, que ha sufrido reestructuraciones y reducciones de plantilla, ¿cómo se intenta gestionar un proceso tan delicado como un despido colectivo (ERE) o salidas individuales para que el impacto humano sea el menor posible?}

\textbf{Respuesta:}\newline
``En colectivos queda acotado por el sindicato. En individuales (disciplinarios), ateniéndose a la ley, si se produce no hay forma de mantener al empleado.''

\subsubsection*{Antes de llegar a reducir plantilla, ¿se suelen explorar alternativas dentro del banco, como reasignaciones a otros departamentos, puestos compartidos o cambios en las políticas de empleo? ¿Creéis que hay margen para incorporar nuevas medidas?}

\textbf{Respuesta:} \newline
``Se intenta en la medida de lo posible, pero no suele ocurrir.''

\subsubsection*{Por último, hemos estudiado que los programas de recolocación u outplacement ayudan mucho a los empleados que salen a encontrar un nuevo rumbo profesional. ¿CaixaBank ofrece este tipo de servicios o apoyo emocional y formativo para facilitar su integración en el mercado laboral? ¿Qué tal es la experiencia con esto?}

\textbf{Respuesta:}\newline
``No se trabaja, pero se ha trabajado en época por ejemplo de la fusión. Proporcionándoles la entrada en empresas del mismo grupo.''

